% OVERLEAF WILL NOT BUILD THIS FILE. It \inputs docs/latex/coursework_preamble.tex by a
% relative path that climbs above the project root (so does reference_docs/cesc470_macros.tex),
% and an Overleaf project is self-contained: "File `../../../../docs/latex/coursework_preamble.tex' not found."
% Upload /home/devel/electrical_notes/content/cesc_470/hw/hw01/overleaf/p10_amdahl_speedup.tex instead -- docs/latex/flatten_tex.sh writes it.
% Edit THIS file, never that generated copy -- docs/latex/flatten_tex.sh overwrites it. Regenerate with:
%   docs/latex/flatten_tex.sh content/cesc_470/hw/hw01   (run from /home/devel/electrical_notes)
% ─────────────────────────────────────────────────────────────────────────────
% SHARED coursework preamble — every course inputs THIS.
%
% PREAMBLE ONLY. No \begin{document}, no course-specific notation.
%
% Course-specific macros live at COURSE level, in
%   content/<course>/reference_docs/<course>_macros.tex
% which is inputted AFTER this file. ONE per course, shared by hw/qz/exam --
% macros under a single kind are unreachable from a sibling kind.
%
% That split is the point: the page style, the answer box, and the problem
% header are identical everywhere, so they are defined once; \conv and \ustep
% mean something in DSP and nothing in a networks course, so they are not.
%
% A per-problem file therefore starts:
%   % ─────────────────────────────────────────────────────────────────────────────
% SHARED coursework preamble — every course inputs THIS.
%
% PREAMBLE ONLY. No \begin{document}, no course-specific notation.
%
% Course-specific macros live at COURSE level, in
%   content/<course>/reference_docs/<course>_macros.tex
% which is inputted AFTER this file. ONE per course, shared by hw/qz/exam --
% macros under a single kind are unreachable from a sibling kind.
%
% That split is the point: the page style, the answer box, and the problem
% header are identical everywhere, so they are defined once; \conv and \ustep
% mean something in DSP and nothing in a networks course, so they are not.
%
% A per-problem file therefore starts:
%   % ─────────────────────────────────────────────────────────────────────────────
% SHARED coursework preamble — every course inputs THIS.
%
% PREAMBLE ONLY. No \begin{document}, no course-specific notation.
%
% Course-specific macros live at COURSE level, in
%   content/<course>/reference_docs/<course>_macros.tex
% which is inputted AFTER this file. ONE per course, shared by hw/qz/exam --
% macros under a single kind are unreachable from a sibling kind.
%
% That split is the point: the page style, the answer box, and the problem
% header are identical everywhere, so they are defined once; \conv and \ustep
% mean something in DSP and nothing in a networks course, so they are not.
%
% A per-problem file therefore starts:
%   \input{../../../../docs/latex/coursework_preamble.tex}
%   \input{../../reference_docs/cesc470_macros.tex}
%   \begin{document}
%
% Build-check one file, or a whole assignment:
%   docs/latex/build_tex.sh content/cesc_470/hw/hw01/p01_five_components.tex
%   docs/latex/build_tex.sh content/cesc_470/hw/hw01
%
% Doctrine: docs/directives/coursework-solutions.md
% ─────────────────────────────────────────────────────────────────────────────

\documentclass[11pt]{article}

\usepackage[margin=1in]{geometry}
\usepackage{amsmath, amssymb, mathtools}
\usepackage{graphicx}
\usepackage{float}
\usepackage{booktabs}
\usepackage{longtable}
\usepackage{enumitem}
\usepackage{siunitx}
\usepackage[dvipsnames]{xcolor}
\usepackage{tcolorbox}
\usepackage{fancyhdr}
\usepackage{caption}
\usepackage{tikz}
\usepackage{pgfplots}
\pgfplotsset{compat=1.18}
\usepackage[colorlinks=true, allcolors=NavyBlue]{hyperref}

\captionsetup{font=small, labelfont=bf}
\setlist[enumerate]{itemsep=2pt, topsep=4pt}
\setlist[itemize]{itemsep=2pt, topsep=4pt}

% --- final-answer box -------------------------------------------------------
% Every problem file ends with its result in one of these. The condensed
% solutions document is assembled by lifting exactly these.
\newtcolorbox{answerbox}[1][]{
  colback=Green!6, colframe=Green!55!black,
  boxrule=0.6pt, arc=2pt, left=6pt, right=6pt, top=4pt, bottom=4pt,
  #1
}

% --- citation to course material --------------------------------------------
% \cite{Module 01, PDF p55}  ->  a small italic parenthetical.
% Solutions cite the SLIDE/PAGE a formula came from so a grader (or a future
% reader) can check the claim against what was actually taught. See the
% directive: every non-obvious step carries one.
\newcommand{\srcref}[1]{{\small\itshape(#1)}}

% --- a step that came from OUTSIDE the course materials ---------------------
% Used only where the course is genuinely silent and the problem asks for
% outside research. Marked visually so it is never mistaken for lecture content.
\newcommand{\extref}[1]{{\small\itshape\color{Sepia}[external: #1]}}

% --- callout for the trap in a problem --------------------------------------
\newtcolorbox{trap}[1][]{
  colback=BurntOrange!7, colframe=BurntOrange!70!black,
  boxrule=0.6pt, arc=2pt, left=6pt, right=6pt, top=4pt, bottom=4pt,
  fonttitle=\bfseries, title=Watch out, #1
}

% --- per-problem header -----------------------------------------------------
% \problemheader{8}{15 pts}{Target clock rate}
% The optional 4-argument form carries a learning outcome where the course
% states one (CESC 410 does; CESC 470's HW1 does not).
\newcommand{\problemheader}[3]{%
  \begin{center}\large\bfseries
    Problem #1 \quad\normalsize\mdseries(#2)\\[2pt]
    \large\bfseries #3
  \end{center}
  \vspace{2pt}\hrule\vspace{8pt}
}
\newcommand{\problemheaderlo}[4]{%
  \begin{center}\large\bfseries
    Problem #1 \quad\normalsize\mdseries(#2, #3)\\[2pt]
    \large\bfseries #4
  \end{center}
  \vspace{2pt}\hrule\vspace{8pt}
}

% Course name in the running head. Defined BEFORE \lhead references it, and
% \renewcommand'd by each course's macros file (inputted after this one).
\providecommand{\coursename}{}

\pagestyle{fancy}
\fancyhf{}
\lhead{\small\coursename}
\rhead{\small\thepage}
\renewcommand{\headrulewidth}{0.3pt}

%   % CESC 470 Computer Architecture — course-specific macros ONLY.
%
% Inputted AFTER docs/latex/coursework_preamble.tex, which carries everything
% shared (answerbox, problemheader, srcref, page style). Put nothing general
% here: if another course would want it, it belongs in the shared preamble.

\renewcommand{\coursename}{CESC 470}

% --- performance-equation notation ------------------------------------------
% The course writes these in words on the slides; these render them compactly
% and, more importantly, IDENTICALLY across every problem file.
\newcommand{\CPUtime}{T_{\mathrm{CPU}}}
\newcommand{\IC}{\mathrm{IC}}                    % instruction count
\newcommand{\CPI}{\mathrm{CPI}}
\newcommand{\cycletime}{t_{\mathrm{cycle}}}
\newcommand{\clockrate}{f_{\mathrm{clk}}}
\newcommand{\cycles}{N_{\mathrm{cycles}}}
\newcommand{\speedup}{S}

% Amdahl's law, in the form the slides state it (Module 01, PDF p86, slide 49).
\newcommand{\amdahl}[3]{#1 + \dfrac{#2}{#3}}

%   \begin{document}
%
% Build-check one file, or a whole assignment:
%   docs/latex/build_tex.sh content/cesc_470/hw/hw01/p01_five_components.tex
%   docs/latex/build_tex.sh content/cesc_470/hw/hw01
%
% Doctrine: docs/directives/coursework-solutions.md
% ─────────────────────────────────────────────────────────────────────────────

\documentclass[11pt]{article}

\usepackage[margin=1in]{geometry}
\usepackage{amsmath, amssymb, mathtools}
\usepackage{graphicx}
\usepackage{float}
\usepackage{booktabs}
\usepackage{longtable}
\usepackage{enumitem}
\usepackage{siunitx}
\usepackage[dvipsnames]{xcolor}
\usepackage{tcolorbox}
\usepackage{fancyhdr}
\usepackage{caption}
\usepackage{tikz}
\usepackage{pgfplots}
\pgfplotsset{compat=1.18}
\usepackage[colorlinks=true, allcolors=NavyBlue]{hyperref}

\captionsetup{font=small, labelfont=bf}
\setlist[enumerate]{itemsep=2pt, topsep=4pt}
\setlist[itemize]{itemsep=2pt, topsep=4pt}

% --- final-answer box -------------------------------------------------------
% Every problem file ends with its result in one of these. The condensed
% solutions document is assembled by lifting exactly these.
\newtcolorbox{answerbox}[1][]{
  colback=Green!6, colframe=Green!55!black,
  boxrule=0.6pt, arc=2pt, left=6pt, right=6pt, top=4pt, bottom=4pt,
  #1
}

% --- citation to course material --------------------------------------------
% \cite{Module 01, PDF p55}  ->  a small italic parenthetical.
% Solutions cite the SLIDE/PAGE a formula came from so a grader (or a future
% reader) can check the claim against what was actually taught. See the
% directive: every non-obvious step carries one.
\newcommand{\srcref}[1]{{\small\itshape(#1)}}

% --- a step that came from OUTSIDE the course materials ---------------------
% Used only where the course is genuinely silent and the problem asks for
% outside research. Marked visually so it is never mistaken for lecture content.
\newcommand{\extref}[1]{{\small\itshape\color{Sepia}[external: #1]}}

% --- callout for the trap in a problem --------------------------------------
\newtcolorbox{trap}[1][]{
  colback=BurntOrange!7, colframe=BurntOrange!70!black,
  boxrule=0.6pt, arc=2pt, left=6pt, right=6pt, top=4pt, bottom=4pt,
  fonttitle=\bfseries, title=Watch out, #1
}

% --- per-problem header -----------------------------------------------------
% \problemheader{8}{15 pts}{Target clock rate}
% The optional 4-argument form carries a learning outcome where the course
% states one (CESC 410 does; CESC 470's HW1 does not).
\newcommand{\problemheader}[3]{%
  \begin{center}\large\bfseries
    Problem #1 \quad\normalsize\mdseries(#2)\\[2pt]
    \large\bfseries #3
  \end{center}
  \vspace{2pt}\hrule\vspace{8pt}
}
\newcommand{\problemheaderlo}[4]{%
  \begin{center}\large\bfseries
    Problem #1 \quad\normalsize\mdseries(#2, #3)\\[2pt]
    \large\bfseries #4
  \end{center}
  \vspace{2pt}\hrule\vspace{8pt}
}

% Course name in the running head. Defined BEFORE \lhead references it, and
% \renewcommand'd by each course's macros file (inputted after this one).
\providecommand{\coursename}{}

\pagestyle{fancy}
\fancyhf{}
\lhead{\small\coursename}
\rhead{\small\thepage}
\renewcommand{\headrulewidth}{0.3pt}

%   % CESC 470 Computer Architecture — course-specific macros ONLY.
%
% Inputted AFTER docs/latex/coursework_preamble.tex, which carries everything
% shared (answerbox, problemheader, srcref, page style). Put nothing general
% here: if another course would want it, it belongs in the shared preamble.

\renewcommand{\coursename}{CESC 470}

% --- performance-equation notation ------------------------------------------
% The course writes these in words on the slides; these render them compactly
% and, more importantly, IDENTICALLY across every problem file.
\newcommand{\CPUtime}{T_{\mathrm{CPU}}}
\newcommand{\IC}{\mathrm{IC}}                    % instruction count
\newcommand{\CPI}{\mathrm{CPI}}
\newcommand{\cycletime}{t_{\mathrm{cycle}}}
\newcommand{\clockrate}{f_{\mathrm{clk}}}
\newcommand{\cycles}{N_{\mathrm{cycles}}}
\newcommand{\speedup}{S}

% Amdahl's law, in the form the slides state it (Module 01, PDF p86, slide 49).
\newcommand{\amdahl}[3]{#1 + \dfrac{#2}{#3}}

%   \begin{document}
%
% Build-check one file, or a whole assignment:
%   docs/latex/build_tex.sh content/cesc_470/hw/hw01/p01_five_components.tex
%   docs/latex/build_tex.sh content/cesc_470/hw/hw01
%
% Doctrine: docs/directives/coursework-solutions.md
% ─────────────────────────────────────────────────────────────────────────────

\documentclass[11pt]{article}

\usepackage[margin=1in]{geometry}
\usepackage{amsmath, amssymb, mathtools}
\usepackage{graphicx}
\usepackage{float}
\usepackage{booktabs}
\usepackage{longtable}
\usepackage{enumitem}
\usepackage{siunitx}
\usepackage[dvipsnames]{xcolor}
\usepackage{tcolorbox}
\usepackage{fancyhdr}
\usepackage{caption}
\usepackage{tikz}
\usepackage{pgfplots}
\pgfplotsset{compat=1.18}
\usepackage[colorlinks=true, allcolors=NavyBlue]{hyperref}

\captionsetup{font=small, labelfont=bf}
\setlist[enumerate]{itemsep=2pt, topsep=4pt}
\setlist[itemize]{itemsep=2pt, topsep=4pt}

% --- final-answer box -------------------------------------------------------
% Every problem file ends with its result in one of these. The condensed
% solutions document is assembled by lifting exactly these.
\newtcolorbox{answerbox}[1][]{
  colback=Green!6, colframe=Green!55!black,
  boxrule=0.6pt, arc=2pt, left=6pt, right=6pt, top=4pt, bottom=4pt,
  #1
}

% --- citation to course material --------------------------------------------
% \cite{Module 01, PDF p55}  ->  a small italic parenthetical.
% Solutions cite the SLIDE/PAGE a formula came from so a grader (or a future
% reader) can check the claim against what was actually taught. See the
% directive: every non-obvious step carries one.
\newcommand{\srcref}[1]{{\small\itshape(#1)}}

% --- a step that came from OUTSIDE the course materials ---------------------
% Used only where the course is genuinely silent and the problem asks for
% outside research. Marked visually so it is never mistaken for lecture content.
\newcommand{\extref}[1]{{\small\itshape\color{Sepia}[external: #1]}}

% --- callout for the trap in a problem --------------------------------------
\newtcolorbox{trap}[1][]{
  colback=BurntOrange!7, colframe=BurntOrange!70!black,
  boxrule=0.6pt, arc=2pt, left=6pt, right=6pt, top=4pt, bottom=4pt,
  fonttitle=\bfseries, title=Watch out, #1
}

% --- per-problem header -----------------------------------------------------
% \problemheader{8}{15 pts}{Target clock rate}
% The optional 4-argument form carries a learning outcome where the course
% states one (CESC 410 does; CESC 470's HW1 does not).
\newcommand{\problemheader}[3]{%
  \begin{center}\large\bfseries
    Problem #1 \quad\normalsize\mdseries(#2)\\[2pt]
    \large\bfseries #3
  \end{center}
  \vspace{2pt}\hrule\vspace{8pt}
}
\newcommand{\problemheaderlo}[4]{%
  \begin{center}\large\bfseries
    Problem #1 \quad\normalsize\mdseries(#2, #3)\\[2pt]
    \large\bfseries #4
  \end{center}
  \vspace{2pt}\hrule\vspace{8pt}
}

% Course name in the running head. Defined BEFORE \lhead references it, and
% \renewcommand'd by each course's macros file (inputted after this one).
\providecommand{\coursename}{}

\pagestyle{fancy}
\fancyhf{}
\lhead{\small\coursename}
\rhead{\small\thepage}
\renewcommand{\headrulewidth}{0.3pt}

% CESC 470 Computer Architecture — course-specific macros ONLY.
%
% Inputted AFTER docs/latex/coursework_preamble.tex, which carries everything
% shared (answerbox, problemheader, srcref, page style). Put nothing general
% here: if another course would want it, it belongs in the shared preamble.

\renewcommand{\coursename}{CESC 470}

% --- performance-equation notation ------------------------------------------
% The course writes these in words on the slides; these render them compactly
% and, more importantly, IDENTICALLY across every problem file.
\newcommand{\CPUtime}{T_{\mathrm{CPU}}}
\newcommand{\IC}{\mathrm{IC}}                    % instruction count
\newcommand{\CPI}{\mathrm{CPI}}
\newcommand{\cycletime}{t_{\mathrm{cycle}}}
\newcommand{\clockrate}{f_{\mathrm{clk}}}
\newcommand{\cycles}{N_{\mathrm{cycles}}}
\newcommand{\speedup}{S}

% Amdahl's law, in the form the slides state it (Module 01, PDF p86, slide 49).
\newcommand{\amdahl}[3]{#1 + \dfrac{#2}{#3}}


\begin{document}

\problemheader{10}{10 pts}{Amdahl's law --- speedup from a $10\times$ FP improvement}

\subsection*{Problem statement}

A program spends 80\% of its time performing floating-point operations. If
floating-point performance is improved by a factor of 10, what is the maximum
possible speedup?

\subsection*{Approach}

Amdahl's law, in the form the course states it \srcref{Module 01, PDF p86,
slide 49}:

\begin{quote}
``Speed up overall $=$ Execution Time Unaffected $+$ ( Execution Time Affected
$/$ Amount of Improvement )''
\end{quote}

which, read as the slide's worked example uses it, gives the \emph{new
execution time}:
\begin{equation}
  \CPUtime^{\text{new}} = \amdahl{T_{\text{unaffected}}}{T_{\text{affected}}}{n},
  \qquad
  \speedup = \frac{\CPUtime^{\text{old}}}{\CPUtime^{\text{new}}}.
  \label{eq:amdahl}
\end{equation}

No absolute run time is given, only a percentage, so normalise the original time
to $1$. Then $T_{\text{affected}} = 0.8$ and $T_{\text{unaffected}} = 0.2$, and
the speedup comes out dimensionless — which it must, being a ratio.

\subsection*{Work}

\paragraph{Step 1 --- split the time.}
Take $\CPUtime^{\text{old}} = 1$ (any unit; the ratio is what matters).
\begin{align}
  T_{\text{affected}}   &= 0.80 \quad\text{(floating point)} \\
  T_{\text{unaffected}} &= 1 - 0.80 = 0.20 \quad\text{(everything else)}
\end{align}

\paragraph{Step 2 --- apply the improvement.}
Only the affected part is divided by $n = 10$; the rest is untouched. That
asymmetry is the entire content of Amdahl's law:
\begin{align}
  \CPUtime^{\text{new}} &= 0.20 + \frac{0.80}{10} \\
                        &= 0.20 + 0.08 \\
                        &= 0.28.
\end{align}

\paragraph{Step 3 --- form the speedup.}
\begin{equation}
  \speedup = \frac{\CPUtime^{\text{old}}}{\CPUtime^{\text{new}}}
           = \frac{1}{0.28}
           = \mathbf{3.571\ldots \approx 3.57{\times}}.
\end{equation}

\paragraph{Step 4 --- the ceiling, which is what ``maximum possible'' probes.}
Even with \emph{infinitely} fast floating point ($n \to \infty$), the $0.20$
never goes away:
\begin{equation}
  \speedup_{\max} = \lim_{n\to\infty} \frac{1}{0.20 + \dfrac{0.80}{n}}
                  = \frac{1}{0.20} = 5{\times}.
\end{equation}
So $3.57\times$ of a possible $5\times$ — a $10\times$ improvement to 80\% of the
work buys about 71\% of the benefit that infinite improvement would.

This mirrors the course's own example exactly \srcref{PDF p87, slide 50}, where
a program of 100 s with 80 s of multiply can be made $4\times$ faster (needing
$n = 16$) but \textbf{cannot} be made $5\times$ faster at any $n$ — the same
$1/0.20 = 5$ ceiling appearing as an impossibility rather than a limit.

\begin{trap}
\textbf{The question is ambiguous, and the two readings give different numbers.}
\begin{itemize}[nosep]
  \item ``What speedup results from the stated $10\times$ improvement?''
    $\Rightarrow \mathbf{3.57\times}$
  \item ``What is the maximum speedup obtainable by improving this part at all?''
    $\Rightarrow \mathbf{5\times}$
\end{itemize}
Since the problem supplies a specific factor of 10, \textbf{3.57$\times$ is the
answer to give}, with the $5\times$ ceiling stated alongside it. Giving only
$5\times$ ignores the data provided; giving only $3.57\times$ misses what
``maximum possible'' is pointing at. Present both.
\end{trap}

\subsection*{Check}

\paragraph{Concrete numbers instead of fractions.}
Suppose the program takes 100 s: 80 s floating point, 20 s other. After a
$10\times$ FP improvement the FP work takes $80/10 = 8$ s, so
\begin{equation}
  \CPUtime^{\text{new}} = 20 + 8 = 28\ \text{s},
  \qquad \speedup = \frac{100}{28} = 3.571 \checkmark
\end{equation}
matching the normalised calculation, and confirming the normalisation was safe.

\paragraph{Bounds.}
$1 < 3.57 < 5$. It must exceed 1 (the change helps) and stay below the $1/0.2 = 5$
ceiling (the unaffected 20\% is irreducible) \checkmark.

\paragraph{Where the remaining time goes.}
After the improvement, the 28 s splits as 20 s unaffected and 8 s floating point
--- so the \emph{unaffected} portion has grown from 20\% to $20/28 \approx 71\%$
of run time. Further FP optimisation would now buy very little, which is the
practical lesson: Amdahl's law is why the course's design principle is
\emph{make the common case fast} \srcref{PDF p86, slide 49} --- and why the
common case shifts once you have optimised it.

\subsection*{Answer}

\begin{answerbox}
With $T_{\text{affected}} = 0.80$, $T_{\text{unaffected}} = 0.20$, $n = 10$:
$$\CPUtime^{\text{new}} = 0.20 + \frac{0.80}{10} = 0.28
\qquad\Longrightarrow\qquad
\speedup = \frac{1}{0.28} = \mathbf{3.57\times}$$

\medskip
\textbf{Theoretical ceiling:} even with infinitely fast floating point, the
untouched 20\% remains, so
$$\speedup_{\max} = \frac{1}{0.20} = \mathbf{5\times}$$

\medskip
{\small The $10\times$ improvement therefore achieves $3.57$ of the $5\times$
that is possible in the limit. Give \textbf{3.57$\times$} as the answer to the
stated improvement, and note the \textbf{5$\times$} ceiling.}
\end{answerbox}

\end{document}
